\par\addvspace{7pt}\Needspace{4\baselineskip}\noindent\textbf{课堂题：同构、染色、最小生成树与匹配}\par
\begin{exercise}\label{mit:6042-cp19-q1}\label{mit:6042-cp32-q1}
\textbf{6.042J，Spring 2015，CP19 第 1 题；同题 CP32 第 1 题。}
一个群体有 $m$ 位男性、$f$ 位女性；每条交往关系恰连接其中一男一女。观察到男性平均女性伴侣数比女性平均男性伴侣数高 $10\%$。
（a）评价“平均数应相同，所以男性必定夸大”的断言；（b）求 $m=cf$ 中的 $c$；（c）若女性 $20\%$、男性 $5\%$ 没有任何伴侣，排除他们后，男性平均数是女性平均数的 $x(f/m)$ 倍，求 $x$；（d）为何不能将每位女性与不同男性配对？
\end{exercise}
\sourceof{CP19 与 CP32，物理第 1 页；两卷正文重复，编者独立解答。}
\begin{solution}
用二分图表示，边数为 $e>0$，两侧的度数和都为 $e$。
(a) 平均数分别是 $e/m,e/f$，分母不同，未必相等，不能据此断定夸大。
(b) $e/m=1.1e/f$，所以 $c=10/11$。
(c) 排除孤立点不改变 $e$，比值为 $(e/(0.95m))/(e/(0.8f))=(16/19)(f/m)$，故 $x=16/19$。
(d) $m<f$，从全部女性到男性不可能有单射，甚至不要求配对是原图的边也不可能。
\end{solution}

\begin{exercise}\label{mit:6042-cp19-q2}\label{mit:6042-cp32-q2}
\textbf{6.042J，Spring 2015，CP19 第 2 题；同题 CP32 第 2 题。}
（a）证明简单图中奇度顶点数为偶数；（b）将其解释为聚会握手结论；（c）若 George 握手次数为奇数，证明从他出发有一条握手序列到达另一位握手次数为奇数的人。
\end{exercise}
\sourceof{CP19 与 CP32，物理第 1 页；编者独立解答。}
\begin{solution}
(a) 每条边对度数和贡献 $2$，故总度数是偶数。偶度项对模 $2$ 的和没有贡献，每个奇度项贡献 $1$，所以奇度项数为偶数。
(b) 人是顶点、发生握手的两人间连一边，度数就是握手次数，直接应用 (a)。
(c) 在 George 所在连通分支中单独应用 (a)。既有一个奇度顶点，该分支必有另外一个奇度顶点 $h$。连通性给 George 到 $h$ 的路径，沿路径的人组成所需序列。这比只在全图寻找另一奇度点更强，因为它保证可达。
\end{solution}

\begin{exercise}\label{mit:6042-cp19-q3}\label{mit:6042-cp32-q3}
\textbf{6.042J，Spring 2015，CP19 第 3 题；同题 CP32 第 3 题。}
左图顶点为 $1,\ldots,6$，边集为 $13,14,23,24,35,56,64$；右图顶点为 $a,b,c,d,e,f$，边集为 $ac,ad,bc,bd,ce,ef,fd$。列出全部同构，并解释为何没有其他同构。
\end{exercise}
\sourceof{CP19 与 CP32，物理第 1--2 页；原图边集重构，编者独立解答。}
\begin{solution}
按 $1,2,3,4,5,6$ 的顺序列像，全部四个同构是
\[
(a,b,c,d,e,f),\quad(b,a,c,d,e,f),\quad
(a,b,d,c,f,e),\quad(b,a,d,c,f,e).
\]
只有 $3,4$ 度为 $3$，须映到 $c,d$，有两个方向。顶点 $1,2$ 都同时邻接这两个度 $3$ 顶点，须映到 $a,b$，有两种顺序；最后从 $3$ 到 $4$ 的长 $3$ 路径 $3,5,6,4$ 的内部顶点像随端点方向唯一确定。故最多四个，上列均保持边，恰是全部。
\end{solution}

\begin{exercise}\label{mit:6042-cp19-q4}\label{mit:6042-cp32-q4}
\textbf{6.042J，Spring 2015，CP19 第 4 题；同题 CP32 第 4 题。}
判定下列性质是否为简单图同构不变量：
\begin{enumerate}[label=(\alph*)]
\item 顶点可编号为 $1$ 至 $7$；\item 有经过全部顶点的圈；\item 有两个度 $8$ 顶点；
\item 当前图画中两条边等长；\item 无论删哪一边，任意两点之间仍有路径；
\item 有两个顶点不交的圈；\item 某顶点作为集合是另一顶点的子集；
\item 图可以画成所有边都等长；\item 两个同构不变量的逻辑“或”；\item 一个同构不变量的否定。
\end{enumerate}
\end{exercise}
\sourceof{CP19，物理第 1--2 页；CP32 相同；编者独立解答。}
\begin{solution}
依次为：是、是、是、否、是、是、否、是、是、是。
(a) 等价于有 $7$ 个顶点；(b) 同构把 Hamilton 圈映成 Hamilton 圈；(c) 邻点对应使度数保持；(e) 删边后的对应仍保持路径；(f) 单射保持两圈不交。
(d) 长度属于某次画法，同一抽象图可以重新摆放顶点改变它；(g) 顶点的名字可以是集合，也可以全改成数字，包含关系不属于边结构。
(h) 虽涉及几何，但它说的是\textbf{存在}一种画法：若 $G$ 有这样的画法，把同构图 $H$ 的对应顶点放在相同位置，边的等长性质就转移过去；反向同理。
(i)(j) 若两个图对每个原性质的真假都相同，那么对这些真值取“或”或否定，结果仍相同。
\end{solution}

\begin{exercise}\label{mit:6042-cp19-q5}\label{mit:6042-cp32-q5}
\textbf{6.042J，Spring 2015，CP19 第 5 题（补充题）；同题 CP32 第 5 题。}
有向图 $G$ 的出邻点集记为 $G(v)$。若 $f:G\to H$ 是图同构，（a）仅用定义证明 $f(G(v))=H(f(v))$；（b）推出对每个 $k\in\N$，两图的出度 $k$ 顶点数相同。
\end{exercise}
\sourceof{CP19 与 CP32，物理第 2 页；编者独立解答。}
\begin{solution}
(a) 对任意 $h\in V(H)$，令 $x=f^{-1}(h)$，则
\[
h\in H(f(v))\iff(f(v),f(x))\in E(H)
\iff(v,x)\in E(G)\iff h\in f(G(v)).
\]
这逐元素证明集合相等。
(b) $f$ 限制在 $G(v)$ 上仍为双射，故 $|G(v)|=|H(f(v))|$。因此 $f$ 也将出度 $k$ 的顶点集双射到另一图相应集合，基数相同。
\end{solution}

\begin{exercise}\label{mit:6042-cp20-q1}
\textbf{6.042J，Spring 2015，CP20 第 1 题。}
输入为 $a,b$；依次计算 $c=a+b$，$d=ac$，$e=c+3$，$f=c-e$，$g=a+f$，$h=f+1$；输出为 $d,g,h$。不再使用的变量寄存器可覆盖，包括最后一次读完该变量的同一步。
（a）将寄存器分配化为图染色并构造冲突图；（b）给最少寄存器分配；（c）若变量重复赋值，例如 $t=r+s,u=3t,t=m-k,v=t+u$，应如何处理？
\end{exercise}
\sourceof{CP20，物理第 1 页；编者独立解答。}
\begin{solution}
(a) 每个仍须保存的变量值为顶点，同时须保留的两个值连边。各步后的活跃集为
$\{a,c\}$、$\{a,c,d\}$、$\{a,c,d,e\}$、$\{a,d,f\}$、$\{d,f,g\}$、$\{d,g,h\}$，初始还有 $\{a,b\}$。所以边集是
$ab,ac,ad,cd,ae,ce,de,af,df,dg,fg,dh,gh$。
(b) 四个寄存器可分为 $R_1:\{a,g\}$，$R_2:\{b,c,f,h\}$，$R_3:\{d\}$，$R_4:\{e\}$，同一组中没有冲突边。$a,c,d,e$ 组成 $K_4$，需要四个不同寄存器，故最少为 $4$。
(c) 将不同赋值产生的值视为不同变量，如 $t_1=r+s,u=3t_1,t_2=m-k,v=t_2+u$，分别确定活跃区间、构造冲突图。最终可把不冲突的 $t_1,t_2$ 分到同一寄存器，但不能事先把它们当作一个必须永久保存的值。
\end{solution}

\begin{exercise}\label{mit:6042-cp20-q2}
\textbf{6.042J，Spring 2015，CP20 第 2 题。}
（a）反驳“每个顶点度数为正的图必连通”；（b）找下列伪证的首个缺口：$n\le2$ 成立；假设 $n$ 顶点结论成立，取一个全正度的 $n$ 顶点图，它连通；增加一个正度新顶点 $x$，它邻接旧顶点 $y$，故通过 $y$ 可达所有旧顶点；于是所有 $n+1$ 顶点图成立。
\end{exercise}
\sourceof{CP20，物理第 2 页；编者独立解答。}
\begin{solution}
(a) 两条互不相交的边 $K_2\sqcup K_2$ 中每点度 $1$，但不连通。
(b) 无依据的是把任意全正度的 $n+1$ 顶点图视为由一个全正度的 $n$ 顶点图增加顶点产生。删去 $x$ 后，其唯一邻点可能变成孤立点，旧图未必满足归纳假设。例如从反例删任何顶点，旧图都会有一个孤立点。证明只说明在\textbf{已经连通}的图上新增一个邻接旧图的顶点保持连通，不能覆盖目标类。
\end{solution}

\begin{exercise}\label{mit:6042-cp20-q3}
\textbf{6.042J，Spring 2015，CP20 第 3 题。}
原图的三色逻辑器件用边集自足表示。$T,F,N$ 为一个三角形，输入为 $P,Q$，输出为 $O$，辅助点为 $u,v,x,z$。其他边是
\[
NP,NQ,NO,\ Tx,xO,xz,\ Ou,Ov,uv,uP,vQ,\ Pz,Qz.
\]
颜色以三角形顶点的颜色命名。（a）证明给定 $P,Q$ 的颜色可扩为三染色，当且仅当二者都不是 $N$；（b）证明 $O$ 为 $T$ 当且仅当 $P,Q$ 至少一个为 $T$；（c）改变一条边的一个端点，使之模拟逻辑“与”。
\end{exercise}
\sourceof{CP20，物理第 2--3 页，图 1；编者独立解答，原图弧线均为无向边。}
\begin{solution}
(a) $NP,NQ$ 立即排除色 $N$。其余四种输入确有扩张，下表给出一组完整赋色：
\[
\begin{array}{cc|ccccc}
P&Q&O&u&v&x&z\\\hline
T&T&T&F&N&F&N\\T&F&T&F&N&F&N\\
F&T&T&N&F&F&N\\F&F&F&T&N&N&T
\end{array}
\]
代入全部边，每条边两端颜色均不同。

(b) $Ouv$ 是三角形。若 $P=Q=T$，则 $u,v$ 都不能为 $T$，所以 $O=T$；若 $P=Q=F$，同理 $O=F$。若 $P,Q$ 异色，$z$ 同时邻接二者，须为 $N$；$x$ 邻接 $z,T$，须为 $F$；$O$ 邻接 $x,N$，须为 $T$。因此任何扩张的输出都实现“或”，而不只是表中选定的扩张。

(c) 把 $Tx$ 改为 $Fx$。输入同色时仍由三角形 $Ouv$ 强制输出同色；输入异色时仍有 $z=N$，现强制 $x=T$、$O=F$。扩张存在性可由表逐项给出：对输入 $TT,TF,FT,FF$，$(O,u,v,x,z)$ 分别取 $(T,F,N,N,F)$、$(F,N,T,T,N)$、$(F,T,N,T,N)$、$(F,T,N,T,N)$。故恰实现“与”。
\end{solution}

\begin{exercise}\label{mit:6042-cp20-q4}
\textbf{6.042J，Spring 2015，CP20 第 4 题。}
$n$ 维超立方体 $H_n$ 的顶点为 $n$ 位二进制串，恰一位不同的两点相邻。
（a）对 $H_3$ 的任意不同 $x,y$，给出三条除端点外顶点不交的路径；（b）求 $\kappa(H_3)$；（c）将论证推广到 $H_4$，并尝试一般 $H_n$。
\end{exercise}
\sourceof{CP20，物理第 3 页；编者独立解答。}
\begin{solution}
(a) 翻转固定位置和置换坐标都保持邻接，因此只需令 $x=000$，按距离讨论三个目标。下列每行列出三条路径：
\[
\begin{array}{c|l}
y=100&000,100;\quad000,010,110,100;\quad000,001,101,100\\
y=110&000,100,110;\quad000,010,110;\quad000,001,011,111,110\\
y=111&000,100,110,111;\quad000,010,011,111;\quad000,001,101,111
\end{array}
\]
相邻串均只差一位，各行内部顶点集两两不交。

(b) 删至多两个顶点，任意两个剩余点的三条路径至少有一条不被内部删点击中，所以仍连通，$\kappa\ge3$。删去一个顶点的三个邻点可使该点与其他剩余点隔离，故 $\kappa=3$。

(c) 更一般证明删至多 $n-1$ 个顶点不能使 $H_n$ 断开。按最后一位分成两份 $H_{n-1}$，对应顶点之间有完美匹配。以 $n=1$ 为基例。若每份删去至多 $n-2$ 点，归纳假设使每份剩余图连通；对应匹配有 $2^{n-1}$ 条边，而删点至多毁坏 $n-1<2^{n-1}$ 条，故还有跨份边。若某份删去 $n-1$ 点，另一份完整，每个残存点都可通过自己的对应点进入完整份，仍连通。$n\ge2$ 时删去某点全部 $n$ 个邻点会隔离它，故 $\kappa(H_n)=n$；$H_1=K_2$ 也按完全图连通度约定有 $\kappa=1$。特别 $H_4$ 的连通度为 $4$。结合第~\ref{ch:connectivity} 章的 Menger 论证，可得到四条内部不交路径。
\end{solution}

\begin{exercise}\label{mit:6042-cp21-q1}
\textbf{6.042J，Spring 2015，CP21 第 1 题。}
$4\times4$ 网格的顶点为 $(i,j)$，$0\le i,j\le3$。水平边 $h_{i,j}=\{(i,j),(i+1,j)\}$ 的权为 $(4i+j)/100$；竖直边 $v_{j,i}=\{(j,i),(j,i+1)\}$ 的权为 $1+(i+4j)/100$。
（a）逐次选择不成圈的最轻边（Kruskal）；（b）从 $(1,2)$ 开始，每次选恰一端在当前树中的最轻边（Prim）；说明二者的割安全性；（c）从 $(0,0),(0,3),(2,3)$ 三棵单点树并行选各自最轻出边，相遇后用一般割法合并；（d）核对三个结果相同。
\end{exercise}
\sourceof{CP21，物理第 1 页；编者独立解答。}
\begin{solution}
(a) 所有水平边都轻于竖直边，而且水平边仅组成四条不交路径，全部先选。按权顺序是 $h_{0,0},h_{0,1},h_{0,2},h_{0,3}$，接着同样顺序的 $h_{1,j}$、$h_{2,j}$。然后选 $v_{0,0},v_{0,1},v_{0,2}$ 连接四行，得 $15$ 条边，已为生成树。每一步最轻可加入边都是某个当前分支与其余顶点间割上的最轻边，定理~\ref{thm:safeedge} 保证安全。

(b) 边的实际选择次序是
\[
\begin{gathered}
h_{0,2},h_{1,2},h_{2,2},v_{0,1},\quad
h_{0,1},h_{1,1},h_{2,1},v_{0,0},\\
h_{0,0},h_{1,0},h_{2,0},v_{0,2},\quad h_{0,3},h_{1,3},h_{2,3}.
\end{gathered}
\]
它们逐次为当前树顶点集割上的最轻边，故每步也安全。

(c) 三棵树第一轮分别选 $h_{0,0},h_{0,3},h_{1,3}$；后两棵在 $(1,3)$ 相遇，将其合并。已选边形成森林。此后把每个森林分支视为单元，按 (a) 的权序加入未选且连接不同分支的水平边，再加 $v_{0,0},v_{0,1},v_{0,2}$。第一轮每棵树所选边均为对应单点割最轻边；这些边可以一起属于 MST：例如已明确构造的 (a) MST 就包含它们。其后逐次用割安全定理即可。

(d) 三法的边集均为全部 $12$ 条水平边及 $v_{0,0},v_{0,1},v_{0,2}$。权和为 $\sum_{r=0}^{11}r/100+1+1.01+1.02=3.69$。这里全部边权互异，PS8 第 1 题也说明结果必唯一。
\end{solution}

\begin{exercise}\label{mit:6042-cp21-q2}
\textbf{6.042J，Spring 2015，CP21 第 2 题。}
对非空简单图证明：是树，当且仅当任意两个不同顶点之间恰有一条路径。单点图也视为树。
\end{exercise}
\sourceof{CP21，物理第 1 页；编者独立解答，排除空图的真空量词情形。}
\begin{solution}
树连通，故至少有一条路径。若有两条不同路径，在它们首次分离后到首次重合之间取两段，便组成圈，与无圈矛盾。反过来路径存在保证连通；若有圈，任选圈上两个不同顶点，沿圈的两个方向得到两条不同路径，违反唯一性。因此连通且无圈，是树。单点情形没有不同顶点对，且它按定义连通无圈，结论仍成立。
\end{solution}

\begin{exercise}\label{mit:6042-cp21-q3}
\textbf{6.042J，Spring 2015，CP21 第 3 题。}
有限连通带权图有唯一全局最轻边 $e$。证明每棵 MST 都含 $e$。
\end{exercise}
\sourceof{CP21，物理第 1 页；编者独立解答。}
\begin{solution}
若 MST $T$ 不含 $e$，加上 $e$ 形成一个圈。任取圈上不同于 $e$ 的边 $f$，都有 $w(f)>w(e)$，删去 $f$ 得到权更小的生成树，矛盾。这里仅需全局最小者唯一，不必全部边权互异。
\end{solution}

\begin{exercise}\label{mit:6042-cp21-q4}
\textbf{6.042J，Spring 2015，CP21 第 4 题。}
原题说图有“宽度 $1$”，意思是存在顶点次序，使每点在它之前的邻点至多一个，实质是宽度\textbf{至多} $1$，不是后文树宽定义。
（a）证明这样的有限图是森林；（b）证明每棵有限树有此性质，并推出森林与此性质等价。
\end{exercise}
\sourceof{CP21，物理第 2 页；编者独立解答。}
\begin{solution}
(a) 若有圈，在圈上选排序位置最后的顶点，它的两个圈邻点都更早，违反条件。因此无圈，是森林。也可归纳删最后顶点，它至多一条边接到此前森林，不能产生圈。
(b) 对树的顶点数归纳。单点成立；至少两点时有叶 $x$，删叶后仍是树，对剩余树取归纳次序，再把 $x$ 放在最后，它只有一个更早邻点。森林逐分支使用此次序再拼接，因为跨分支没有边。空森林用空次序成立。这也提供期末第 2 题所需的完整归纳证明。
\end{solution}

\begin{exercise}\label{mit:6042-cp22-q1}
\textbf{6.042J，Spring 2015，CP22 第 1 题。}
四位学生和四家公司各按从优到劣排序：
\[
\begin{array}{c|llll}
Albert&HP&Bellcore&AT\&T&Draper\\
Sarah&AT\&T&Bellcore&Draper&HP\\
Tasha&HP&Draper&AT\&T&Bellcore\\
Elizabeth&Draper&AT\&T&Bellcore&HP\\\hline
AT\&T&Elizabeth&Albert&Tasha&Sarah\\
Bellcore&Tasha&Sarah&Albert&Elizabeth\\
HP&Elizabeth&Tasha&Albert&Sarah\\
Draper&Sarah&Elizabeth&Tasha&Albert
\end{array}
\]
（a）用延迟接受算法求两个稳定分配；（b）给判定任意此类问题的稳定匹配是否唯一的程序。
\end{exercise}
\sourceof{CP22，物理第 1 页；编者独立解答。}
\begin{solution}
(a) 学生申请：首轮 Albert、Tasha 申请 HP，Sarah 申请 AT\&T，Elizabeth 申请 Draper。HP 留 Tasha、拒 Albert；Albert 下一次申请 Bellcore，被接受。结果为
$Albert\!:\!Bellcore,\ Sarah\!:\!AT\&T,\ Tasha\!:\!HP,\ Elizabeth\!:\!Draper$。

公司申请：首轮 AT\&T、HP 申请 Elizabeth，Bellcore 申请 Tasha，Draper 申请 Sarah。Elizabeth 留 AT\&T、拒 HP；HP 改申请 Tasha，Tasha 拒 Bellcore；Bellcore 改申请 Sarah，Sarah 拒 Draper；Draper 改申请 Elizabeth，Elizabeth 拒 AT\&T；AT\&T 改申请 Albert。结果为
$Albert\!:\!AT\&T,\ Sarah\!:\!Bellcore,\ Tasha\!:\!HP,\ Elizabeth\!:\!Draper$。
引理~\ref{lem:6042-stable} 保证这两次算法输出都稳定。
(b) 各运行一次两侧申请算法，比较全部配对。相同则唯一，不同则不唯一；正确性正是该引理的极端夹逼结论。本题两个输出不同，所以至少有两个稳定匹配。
\end{solution}

\begin{exercise}\label{mit:6042-cp22-q2}
\textbf{6.042J，Spring 2015，CP22 第 2 题。}
等人数延迟接受算法中，Harry 为男方、Alice 为女方。判定哪些谓词是保持不变量并说明原因：
（a）Alice 是 Harry 名单上唯一的人；（b）存在没有任何申请者的女方；（c）若 Alice 已不在 Harry 名单上，她保留着一位更喜欢于 Harry 的申请者；（d）Alice 已划去且 Harry 比起当前申请对象更喜欢 Alice；（e）若 Alice 在 Harry 名单上，她比起所有当前申请者更喜欢 Harry。
\end{exercise}
\sourceof{CP22，物理第 1 页；编者独立解答。}
\begin{solution}
(a)(c)(d) 是，(b)(e) 不是。
(a) 名单只删不增；若最后一个名字也被划去便名单耗尽，这在等人数完整偏好情况下已由引理~\ref{lem:6042-stable} 排除。
(c) 女方拒绝 Harry 时保留更喜欢的人，之后只会保留更好的申请者。
(d) 已划名字不回来，男方后续申请只沿排序向后，所以此比较保持。
(b) 初始可能有空缺女方，最终全部有伴侣，故该谓词可从真变假。
(e) 可从真变假。具体取三男 Harry、$k,c$，三女 Alice、$X,Y$。Harry 排 $X,Alice,Y$；$k$ 排 $Alice,Y,X$；$c$ 排 $X,Alice,Y$；$X$ 排 Harry、$c,k$，Alice 排 $c$、Harry、$k$。某日 Harry 与 $c$ 都申请 $X$，$k$ 申请 Alice；Alice 更喜欢 Harry 而非当前的 $k$，且仍在 Harry 名单上，所以谓词为真。$X$ 拒 $c$ 后，下一日 $c$ 到 Alice，她改留 $c$，因更喜欢 $c$ 胜过 Harry，谓词变假。$Y$ 的严格排序任意补全。
\end{solution}

\begin{exercise}\label{mit:6042-cp22-q3}
\textbf{6.042J，Spring 2015，CP22 第 3 题。}
已知保持不变量：女方 $G$ 若从男方 $B$ 的名单划去，则她保留着一个比 $B$ 更喜欢的申请者。用它证明算法产生稳定匹配。
\end{exercise}
\sourceof{CP22，物理第 1--2 页；编者独立解答。}
\begin{solution}
终止及完美性由引理~\ref{lem:6042-stable} 的计数论证保证。若终局有阻塞对 $(B,G)$，$B$ 比起终局伴侣更喜欢 $G$，而他的申请按名单从前到后进行，所以此前 $G$ 必已被划去。给定不变量使终局 $G$ 保留的伴侣比 $B$ 更合意，与阻塞条件矛盾。因此终局稳定。
\end{solution}

\begin{exercise}\label{mit:6042-cp22-q4}
\textbf{6.042J，Spring 2015，CP22 第 4 题。}
医院住院岗位分配中，每家医院有不同容量，岗位总数未必等于毕业生人数。如何修改稳定性定义并推广等人数算法？
\end{exercise}
\sourceof{CP22，物理第 2 页；编者独立解答，假设严格偏好及医院按个人排序择优。}
\begin{solution}
学生至多去一家医院，医院 $H$ 至多接收 $q_H$ 位学生。允许未分配和空岗位，将它们置于所有可接受对象之后。阻塞对为学生 $s$ 和医院 $H$：$s$ 未分配或更喜欢 $H$，且 $H$ 有可接受他的空岗位，或比起某位现任学生更喜欢 $s$。

令学生逐次向未拒绝他的最高偏好医院申请。每家医院在新旧申请者中保留至多 $q_H$ 位最喜欢者，拒绝其余人；每个学生对每家医院至多申请一次，故终止。若终局有阻塞对，学生曾先向该医院申请并被拒，此时医院已满且保留的 $q_H$ 人都更合意；以后这 $q_H$ 个名额的质量只会提高，不可能终局有空可接受岗位或较差学生，矛盾。因此稳定。

也可把每家医院拆成 $q_H$ 个同偏好岗位，学生把同院岗位连续排列并用固定次序打破并列，再加虚拟学生或虚拟岗位平衡人数。虚拟对象排在所有真实可接受对象之后，把结果中含虚拟对象的配对删去。若医院模型中出现阻塞，学生可与该院一个空缺或较差的真实岗位组成岗位模型阻塞，故上述转换正确。若允许不可接受对象，则双方把它们排在虚拟对象之后，终局只保留可接受的真实配对。
\end{solution}
